% Section 1.4, from lectures/L01/appendix/c_exercises.md.

\section{Exercises}
\label{c1:sec:exercises}\label{c1:app:c}

\begin{aside}
\textbf{How to work these.} These exercises practice the concepts from \secref{c1:sec:cpp}. Each set
is one program: write it in its own directory under \file{lectures/L01/exercises} in the companion
repository, with the file names its README gives, and build and run it with the Makefile from
\secref{c1:sec:makefile}.

\textbf{How to check your work.} Run \code{make test} from the repository root. A set's tests
switch on as soon as its files exist, and compare your program's output with the example output
each exercise gives, character for character; a set you have not started is reported as
\code{SKIP}.

The solutions are published in the companion repository, one program per set under
\file{lectures/L01/appendix/solutions}, each with its Makefile. Try each exercise yourself before
looking at them.
\end{aside}

% ------------------------------------------------------------------------------------------
\exerciseset{Default Arguments and Functions}
\label{c1:set:1}

\exercise{Debug logger}{Code}
\label{c1:ex:1.1}
Create a namespace called \code{debug} containing a function:

\begin{cppcode}
void log(const char* message, const std::uint8_t level = 0U);
\end{cppcode}

\task{Tasks}
\begin{enumerate}
  \item Implement the function so it prints the message using \code{std::printf} from
    \code{<cstdio>}.
  \item Include the log level in the output.
  \item Mark the function \code{noexcept}.
\end{enumerate}
Example usage:

\begin{cppcode}
debug::log("System started");
debug::log("Sensor failure", 2U);
\end{cppcode}

\noindent Example output:

\begin{console}
System started, log level = 0
Sensor failure, log level = 2
\end{console}

\exercise{System delay}{Code}
\label{c1:ex:1.2}
In the same program as in Exercise~\ref{c1:ex:1.1}, create a function that simulates a simple
software delay:

\begin{cppcode}
void delay_ms(const std::uint32_t ms = 1U);
\end{cppcode}

\task{Tasks}
\begin{enumerate}
  \item Place the function inside namespace \code{app}.
  \item Mark the function \code{noexcept}.
  \item Inside the function, define a \code{constexpr} constant named \code{maxCount} and assign it
    the value \code{10000000UL}.
  \item Implement a delay using nested loops:
  \begin{itemize}
    \item The outer loop should iterate \code{ms} times.
    \item The inner loop should iterate up to \code{maxCount}.
    \item To reduce the risk of the compiler optimizing the loop away, add a volatile dummy variable
      \code{volatile std::uint32_t dummy{}} inside the function and increment it inside the inner
      loop.
  \end{itemize}
  \item Use this function to simulate a software delay of \code{100 ms} between the debug prints in
    Exercise~\ref{c1:ex:1.1}.
\end{enumerate}

% ------------------------------------------------------------------------------------------
\exerciseset{Struct Driver}
\label{c1:set:2}

\exercise{Software timer}{Code}
\label{c1:ex:2.1}
In this exercise you will implement a simple software timer driver \code{driver::Timer} in a header
file \file{driver/timer.hpp}.

Add the following lines at the top of the file:

\begin{cppcode}
/**
 * @file Timer driver implementation.
 */
#pragma once

#include <cstdint>
#include <cstdio>
\end{cppcode}

\noindent The driver shall simulate a timer that counts milliseconds and generates a timeout when a
configured timeout value is reached.

The struct shall:
\begin{itemize}
  \item Use private member variables with the prefix \code{my}.
  \item Have a constructor.
  \item Have a destructor.
\end{itemize}
Use \code{std::printf} from \code{<cstdio>} for terminal prints.

\task{a) Private member variables}
Add three private member variables:
\begin{itemize}
  \item The first member variable shall:
  \begin{itemize}
    \item Store the timeout in milliseconds.
    \item Have the type \code{const std::uint16_t}.
    \item Be named \code{myTimeout_ms}.
  \end{itemize}
  \item The second member variable shall:
  \begin{itemize}
    \item Store the internal counter in milliseconds.
    \item Have the type \code{std::uint16_t}.
    \item Be named \code{myCounter_ms}.
  \end{itemize}
  \item The third member variable shall:
  \begin{itemize}
    \item Indicate whether the timer is running.
    \item Have the type \code{bool}.
    \item Be named \code{myRunning}.
  \end{itemize}
\end{itemize}

\task{b) Constructor}
Add a constructor that:
\begin{itemize}
  \item Takes two input arguments:
  \begin{itemize}
    \item The timeout in milliseconds.
    \item The initial running state, which shall be set to \code{false} by default.
  \end{itemize}
  \item Initializes:
  \begin{itemize}
    \item \code{myTimeout_ms} with the provided timeout value.
    \item \code{myCounter_ms} with \code{0U}.
    \item \code{myRunning} with \code{false}.
  \end{itemize}
  \item Prints that the timer is being created.
  \item Calls the public method \code{start()} if the initial running state is \code{true}.
\end{itemize}
Example:

\begin{console}
Creating timer!
\end{console}

\task{c) Destructor}
Add a destructor that:
\begin{itemize}
  \item Calls the public method \code{stop()}.
  \item Prints that the timer is being destroyed.
\end{itemize}
Example:

\begin{console}
Destroying timer!
\end{console}

\task{d) Public methods}
Add the following public methods:
\begin{itemize}
  \item \code{timeout_ms()} shall:
  \begin{itemize}
    \item Return \code{myTimeout_ms}.
    \item Be marked \code{const}, \code{noexcept}, and \code{[[nodiscard]]}.
  \end{itemize}
  \item \code{isRunning()} shall:
  \begin{itemize}
    \item Return \code{myRunning}.
    \item Be marked \code{const}, \code{noexcept}, and \code{[[nodiscard]]}.
  \end{itemize}
  \item \code{start()} shall:
  \begin{itemize}
    \item Set \code{myRunning} to \code{true}.
    \item Print \code{Starting timer!}.
  \end{itemize}
  \item \code{stop()} shall:
  \begin{itemize}
    \item Set \code{myRunning} to \code{false}.
    \item Print \code{Stopping timer!}.
  \end{itemize}
  \item \code{toggle()} shall:
  \begin{itemize}
    \item Toggle \code{myRunning}.
    \item Print that the timer was toggled and whether it is now running or stopped:
      \code{Toggling timer: running!} when enabled, \code{Toggling timer: stopped!} when disabled.
  \end{itemize}
  \item \code{tick()} shall:
  \begin{itemize}
    \item Increment \code{myCounter_ms} if \code{myRunning} is \code{true}.
    \item Otherwise do nothing.
  \end{itemize}
  \item \code{hasTimedOut()} shall:
  \begin{itemize}
    \item Return \code{true} if \code{myCounter_ms >= myTimeout_ms}, otherwise \code{false}.
    \item If \code{myCounter_ms >= myTimeout_ms}, reset \code{myCounter_ms} to \code{0U}.
    \item Be marked \code{noexcept} and \code{[[nodiscard]]}.
  \end{itemize}
\end{itemize}

\task{e) Create and use a timer instance}
In \code{main()}:
\begin{itemize}
  \item Create a timer instance:
  \begin{itemize}
    \item Name the instance \code{timer}.
    \item Set the timeout to \code{1000 ms}.
    \item Set the timer to running at startup.
  \end{itemize}
  \item Create a loop that runs for \code{5000} iterations:
  \begin{itemize}
    \item Call \code{tick()} each iteration.
    \item Call \code{hasTimedOut()} to check whether the timer has expired.
    \item Print \code{Timeout after x ms!}, where \code{x} is the configured timeout, whenever the
      timer times out.
  \end{itemize}
\end{itemize}
Since the timer timeout is \code{1000 ms} and the loop runs for \code{5000} iterations, the timer
should generate five timeouts.

Example output:

\begin{console}
Creating timer!
Starting timer!
Timeout after 1000 ms!
Timeout after 1000 ms!
Timeout after 1000 ms!
Timeout after 1000 ms!
Timeout after 1000 ms!
Stopping timer!
Destroying timer!
\end{console}

% ------------------------------------------------------------------------------------------
\exerciseset{References}
\label{c1:set:3}

\exercise{Swap values}{Code}
\label{c1:ex:3.1}
Implement the following function in an anonymous namespace:

\begin{cppcode}
constexpr void swap(std::uint32_t& a, std::uint32_t& b) noexcept;
\end{cppcode}

\task{Tasks}
\begin{enumerate}
  \item Swap the values using a temporary variable \code{temp}.
  \item Test the function with two variables.
\end{enumerate}
Example output:

\begin{console}
Before swap: a = 3, b = 10
After swap: a = 10, b = 3
\end{console}

% ------------------------------------------------------------------------------------------
\exerciseset{Bit Manipulation Templates}
\label{c1:set:4}

\exercise{Clear bit}{Code}
\label{c1:ex:4.1}
Create a function template to clear a bit in a register:

\begin{cppcode}
template<typename T>
constexpr void clear(T& reg, std::uint8_t bit) noexcept;
\end{cppcode}

\noindent Implement this function template inside an anonymous namespace.

\task{Tasks}
\begin{enumerate}
  \item Use \code{static_assert} with \code{std::is_integral<T>::value} from
    \code{<type_traits>} to ensure that \code{T} is an integral type.
  \item Use the error message \code{Cannot perform bit operation with non-integral type!} in the
    \code{static_assert}.
  \item Clear the selected bit in the register.
  \item Print the result in binary using C++ I/O functionality:
  \begin{itemize}
    \item Output stream \code{std::cout} from \code{<iostream>} for printing.
    \item \code{std::bitset<N>} from \code{<bitset>} to generate a bitset representation.
  \end{itemize}
\end{enumerate}
Example usage:

\begin{cppcode}
std::uint8_t reg{0xFFU};
clear(reg, 2U);
std::cout << "Register content: " << std::bitset<8U>(reg) << "\n";
\end{cppcode}

\noindent Expected output:

\begin{console}
Register content: 11111011
\end{console}

\exercise{Toggle bit(s)}{Code}
\label{c1:ex:4.2}
Create a function template for toggling one or several bits in a register using a parameter pack:

\begin{cppcode}
template<typename T, typename... Bits>
constexpr void toggle(T& reg, const Bits... bits) noexcept;
\end{cppcode}

\noindent Implement this function template in the same anonymous namespace as
Exercise~\ref{c1:ex:4.1}.

\task{Tasks}
\begin{enumerate}
  \item Use a parameter pack.
  \item Ensure the type is integral as in Exercise~\ref{c1:ex:4.1}.
  \item Iterate over the bits (for example using \code{{bits...}}).
  \item Toggle all specified bits in the register.
  \item Print the result in binary using C++ I/O functionality:
  \begin{itemize}
    \item Output stream \code{std::cout} from \code{<iostream>} for printing.
    \item \code{std::bitset<N>} from \code{<bitset>} to generate a bitset representation.
  \end{itemize}
\end{enumerate}
Example usage:

\begin{cppcode}
std::uint8_t reg{0xFFU};
clear(reg, 2U);
toggle(reg, 0U, 2U, 4U, 6U);
std::cout << "Register content: " << std::bitset<8U>(reg) << "\n";
\end{cppcode}

\noindent Expected output:

\begin{console}
Register content: 10101110
\end{console}
